\section{Ablation}
\label{sec:ablation}

The protocol has components whose contribution can be removed one at a time: the identity families, the relative tolerance, the separate overlay line, the region qualifier of the non-linearity statement and the weighting of scenarios.
Each removal is measured on the experiment that exercises the component.

\subsection{Identity families}
\label{sec:ablation-families}

\Cref{tab:ablation} leaves out one of the six identity families and reports the share of the \eoneCases{} cases of each defect class, planted at ten times the limit, that the remaining identities still flag.
Exactly \eoneFamiliesLost{} families are irreplaceable for detection: the gross-amount roll-forward, without which a gross-amount movement is never flagged, and the tie of the stage closing to the stage ECL and the overlays, without which a reallocation between stages and a change of the stage ECL are never flagged.
The other four families, the allowance roll-forward and the three waterfall identities, are redundant for single-entry detection, because every class they see is also seen by another family; they remain useful for localization, since they change the number of failing identities from one or two to the values of Section~\ref{sec:res-detection}.
The comparator that keeps only the two roll-forwards of each stage is the removal of the four cross families at once, and it loses four classes.
This is the prediction P3, which held.

\begin{table}[htbp]
\centering
\small
\caption{Ablation of the twelve identities by family: share of the 600 cases of each defect class (defect of ten times the limit) that the remaining identities still flag when one family is left out; the last column counts the classes that no remaining identity detects.}
\label{tab:ablation}
\begin{adjustbox}{max width=\textwidth}
\begin{tabular}{p{4.6cm}rrrrrrrr}
\toprule
Identity family left out & Alw. mov. & Alw. clos. & Realloc. & GCA mov. & WF line & WF clos. & Stage ECL & Classes lost \\
\midrule
Allowance roll-forward by stage & all & all & all & all & all & all & all & 0 \\
Gross amount roll-forward by stage & all & all & all & 0 & all & all & all & 1 \\
Waterfall closes & all & all & all & all & all & all & all & 0 \\
Waterfall equals stage closings & all & all & all & all & all & all & all & 0 \\
Waterfall lines equal stage movements & all & all & all & all & all & all & all & 0 \\
Stage closing equals ECL plus overlays & all & all & 0 & all & all & all & 0 & 2 \\
None left out (twelve identities) & all & all & all & all & all & all & all & 0 \\
\bottomrule
\end{tabular}
\end{adjustbox}
\end{table}


\subsection{The other components}
\label{sec:ablation-components}

\Cref{tab:ablation-components} measures the other components.
Replacing the relative tolerance by an absolute tolerance of $10^{-6}$ units moves the false alarms on clean resampled books from none to \etwoAbsResPct{}.
Removing the cross identities lowers the detection above the limit from \eoneAboveFullK{} of \eoneAboveFullN{} to \eoneAboveStageK{}.
Replacing the separate overlay line by a merge leaves nothing to test: every table closes and the model ECL is overstated by the overlay share over one minus the share.
Dropping the region qualifier makes the statement that the weighted PD is at least the PD at the mean factor false in \efourConcaveNeg{} concave draws and \efourMixedNeg{} mixed draws, against none of the \efourConvexN{} convex draws.
Replacing weighted scenarios by the base scenario alone raises the median absolute error of the expected PD from \efourGhThreeMed{} for the three-point rule to \efourSingleMed{}.

\begin{table}[htbp]
\centering
\small
\caption{Ablation of the other components of the protocol, each measured on the experiment that exercises it: the tolerance basis, the cross identities, the separate overlay line, the region qualifier of the non-linearity statement and the weighting of scenarios, with the value of the measure with and without the component.}
\label{tab:ablation-components}
\begin{adjustbox}{max width=\textwidth}
\begin{tabular}{p{4.8cm}p{3.6cm}p{3.4cm}p{3.4cm}}
\toprule
Component removed & Measure & With the component & Without it \\
\midrule
Relative tolerance basis replaced by an absolute tolerance of 1e-6 units & false alarms on clean books & 0/1,200 [0.0, 0.3] & 68/1,200 [4.5, 7.1] \\
Cross identities removed (stage arithmetic only) & detection above the limit, seven classes & 4,200/4,200 [99.9, 100.0] & 1,800/4,200 [41.4, 44.4] \\
Separate overlay line replaced by a merge into the model ECL & wrong bookings of the overlay detected & 72/72 (four kinds) & not applicable: no line to test \\
Convex-region qualifier dropped from the non-linearity statement & draws that contradict the statement & 0/16,906 (qualified) & 1,844/20,000 (unqualified) \\
Probability-weighted scenarios replaced by the base scenario alone & median absolute error of the expected PD & 0.0030\% (three-point) & 9.20\% \\
\bottomrule
\end{tabular}
\end{adjustbox}
\end{table}

